\documentclass[11pt]{amsart}
\usepackage{geometry} % See geometry.pdf to learn the layout options. There are lots.
\geometry{letterpaper} % ... or a4paper or a5paper or ...
%\geometry{landscape} % Activate for for rotated page geometry
%\usepackage[parfill]{parskip} % Activate to begin paragraphs with an empty line rather than an indent
\usepackage{graphicx}
\usepackage{amssymb}
\usepackage{epstopdf}
\usepackage{multirow}

\DeclareGraphicsRule{.tif}{png}{.png}{convert #1 dirname #1/basename #1 .tif`.png}

\title{Problema 2.4}
%\date{}
% Activate to display a given date or no date

\begin{document}
\maketitle
%\section{}
%\subsection{}
\noindent Construya un diagrama de flujo tal que dado como datos la matrícula y 5 calificaciones de un alumno; imprima la matrícula, el promedio y la palabra “aprobado” si el alumno tiene un promedio mayor o igual que 6, y la palabra “no aprobado” en caso contrario.

\noindent Datos: \hspace{2cm} MAT, CALI, CAL2, CAL3, CAL4, CAL5

\noindent Donde:
\begin{itemize}
    \item \(MAT\) es una variable entera que representa la matrícula del alumno.    
    \item \(CALI, CAL2, CAL3,CAL4 y CAL5\) son variables de tipo real que representan las 5 calificaciones del alumno.
\end{itemize}

\noindent Explicación de las variables
\begin{itemize}
    \item \(MAT\) Variable de tipo entero.
    \item \(CALI, CAL2, CAL3,CAL4 y CAL5\) Variables de tipo real.
    \item \(PRO\) Variable de tipo real. Almacena el promedio de las 5 calificaciones.
\end{itemize}

\noindent A continuación en la siguiente tabla, se puede observar el seguimiento del algoritmo para diferentes corridas.

\begin{tabular}{|c|c|c|c|c|c|c|c|c|}
\hline
\multicolumn{9}{|c|}{\textbf{Tabla 2.11}} \\
\hline
\multicolumn{1}{|c|}{\textbf{NUMERO DE CORRIDA}} & \multicolumn{6}{|c|}{\textbf{DATOS}} & \multicolumn{2}{|c|}{\textbf{RESULTADO}}\\
\hline
 & \textbf{MAT} & \textbf{CAL1} & \textbf{CAL2}  & \textbf{CAL3}  & \textbf{CAL4}  & \textbf{CAL5}  & \textbf{PRO}  & \textbf{COMEMTARIO} \\
\hline
 1 & 16500 & 6 & 7.50 & 8 & 9.50 & 7 & 7.60 & \text{"APROBADO"} \\
\hline
 2 & 16630 & 5 & 4.80 & 7 & 6.30 & 5.90 & 5.80 & \text{"NO APROBADO"} \\
\hline
 3 & 17220 & 8.60 & 9 & 9 & 5.90 & 6.30 & 7.76 & \text{"APROBADO"} \\
\hline
 4 & 18240 & 7 & 4.60 & 4.90 & 7 & 5.60 & 5.82 & \text{"NO APROBADO"} \\
\hline
 5 & 17246 & 8 & 8.50 & 8.30 & 9.20 & 9.30 & 8.66 & \text{"APROBADO"} \\
\hline
 6 & 18250 & 9 & 9.25 & 8.10 & 9.80 & 10 & 9.23 & \text{"APROBADO"} \\
\hline
\end{tabular}







\end{document}